% =============================================================================
% Chapter: Evaluation Framework (system-side). Describes WHICH tests are run and
% WHY; the Results chapter (results_system_draft.tex) presents the numbers.
% =============================================================================

\chapter{Evaluation Framework}
\label{ch:eval-framework}

% -----------------------------------------------------------------------------
\section{Approach}
\label{sec:ef-approach}

This chapter sets out how the \emph{system} is evaluated --- which tests are
applied and why --- while the numbers themselves are reported in
Chapter~\ref{ch:results-system}. It is concerned with operational maturity, not
forecast accuracy: the modelling-side evaluation (MAE, RMSE, sMAPE, rMAE, and the
Diebold--Mariano and Giacomini--White significance tests) is defined in
Chapter~\ref{ch:methods-models} and reported in Chapter~\ref{ch:results-models},
so the two threads of the thesis each carry their own evaluation. Following the
single-case study methodology of Runeson and H\"ost~\cite{runeson2009casestudy},
the platform is evaluated as one system in its real operating context; there is
no controlled comparison against a second platform. The honesty boundary is
stated up front: the figures reported are indicative of the system over the
study window, not claims about long-run production operation. A distinction that
runs through the chapter, and that matters for the validity discussion, is
between \emph{objective measurements} --- numbers produced by reproducible scripts
--- and a single \emph{structured self-assessment}, the ML Test Score
(Table~\ref{tab:ef-test-types}).

\begin{table}[h]
  \centering
  \caption{Two kinds of test used in this work.}
  \label{tab:ef-test-types}
  \begin{tabularx}{\linewidth}{@{}l X X@{}}
    \toprule
    \textbf{Type} & \textbf{What it produces} & \textbf{Auditable how?} \\
    \midrule
    Objective measurement   & A number computed by a script over code, git history, or the database & Re-run the script — same inputs, same number \\
    Structured self-assessment & A rubric score with one-line evidence per item & Read the evidence column; agree or contest each item \\
    \bottomrule
  \end{tabularx}
\end{table}

% -----------------------------------------------------------------------------
\section{Test inventory}
\label{sec:ef-inventory}

Five tests are applied to the system: four objective measurements and one
structured self-assessment, listed in Table~\ref{tab:ef-inventory}. The
remainder of the chapter takes each in turn, stating its definition, the
instrument that produces it, and the reason it is included.

\begin{table}[h]
  \centering
  \caption{Tests applied to the system in this work.}
  \label{tab:ef-inventory}
  \begin{tabularx}{\linewidth}{@{}l l X l@{}}
    \toprule
    \textbf{Test} & \textbf{Type} & \textbf{Source / instrument} & \textbf{Cite} \\
    \midrule
    ML Test Score              & Self-assessment & 28-item rubric, scored 0 / 0.5 / 1, min-of-sections & \cite{breck2017mltestscore} \\
    Provenance coverage        & Objective       & SQL aggregate over \texttt{algo.experiment\_runs} & \cite{gundersen2018reproducibility} \\
    Re-run demonstration       & Objective       & \texttt{scripts/report/reproduce\_run.py} — re-execute logged run, compare to stored prediction & \cite{gundersen2018reproducibility} \\
    DORA delivery throughput   & Objective       & \texttt{gh}/git mining of \texttt{cd\_deploy.yml} history & \cite{forsgren2018accelerate} \\
    Drift monitoring counts    & Objective       & SQL aggregate over \texttt{algo.drift\_reports} & --- \\
    \bottomrule
  \end{tabularx}
\end{table}

% -----------------------------------------------------------------------------
\section{Why these five tests}
\label{sec:ef-why}

Zarour et al.'s systematic review~\cite{zarour2025} catalogues nine MLOps
practices but produces no single quantitative aggregate against which a platform
can be scored. Rather than invent one, this work composes a small set of tests
that each address a measurable face of one or more practices. Between them, the
five tests evidence eight of Zarour's nine practices, as mapped in
Table~\ref{tab:ef-mapping}; the ninth, collaboration and communication, is
qualitative by nature and is documented in the System Implementation chapter
(Chapter~\ref{ch:system-implementation}) rather than measured here.

\begin{table}[h]
  \centering
  \caption{Mapping of tests to the nine MLOps practices in Zarour et al.~\cite{zarour2025}.}
  \label{tab:ef-mapping}
  \begin{tabularx}{\linewidth}{@{}l X@{}}
    \toprule
    \textbf{Zarour practice} & \textbf{Tested by} \\
    \midrule
    P1 Automating CI/CD                  & ML Test Score (Infra), DORA throughput \\
    P2 Orchestrating workflows           & ML Test Score (Infra) \\
    P3 Ensuring reproducibility          & ML Test Score (Infra), Provenance coverage, Re-run demonstration \\
    P4 Versioning                        & ML Test Score (Data, Infra), Provenance coverage \\
    P5 Collaboration and communication   & \emph{not measured here — qualitative, see Ch.~\ref{ch:system-implementation}} \\
    P6 Continuously train and evaluate   & ML Test Score (Model) \\
    P7 Track and log ML metadata         & ML Test Score (Model), Provenance coverage \\
    P8 Continuous monitoring             & ML Test Score (Monitoring), Drift counts \\
    P9 Implementing feedback loops       & Drift counts (re-framed as monitoring; see \S\ref{sec:res-drift}) \\
    \bottomrule
  \end{tabularx}
\end{table}

% -----------------------------------------------------------------------------
\section{ML Test Score}
\label{sec:ef-mlts}

The ML Test Score of Breck et al.~\cite{breck2017mltestscore} is a rubric of 28
tests across four sections --- data, model development, ML infrastructure, and
monitoring --- each scored on the scale $\{0, 0.5, 1\}$ defined in
Table~\ref{tab:ef-mlts-rules}. The final score is the \emph{minimum} of the four
section sums, not their total: because the four sections guard against different
and non-substitutable failure modes (Table~\ref{tab:ef-mlts-sections}), the
weakest section bounds the platform's production-readiness, and the rule rewards
balance rather than breadth. The assessment is conducted as a structured
self-assessment: every item is walked against the codebase, one line of evidence
is recorded for it, and a score is assigned by the explicit rule. The full
per-item evidence is committed to \texttt{Report/metrics/ml\_test\_score.md} so
that each score can be checked or contested independently.

\begin{table}[h]
  \centering
  \caption{Scoring rule applied to each of the 28 items.}
  \label{tab:ef-mlts-rules}
  \begin{tabularx}{\linewidth}{@{}c X@{}}
    \toprule
    \textbf{Score} & \textbf{Meaning} \\
    \midrule
    1.0 & Capability exists and is automated / tested / repeatable (runs in CI, runs on schedule, has a test) \\
    0.5 & Capability exists but is manual, partial, or ad-hoc \\
    0.0 & Not done, or genuinely not applicable to the domain \\
    \bottomrule
  \end{tabularx}
\end{table}

\begin{table}[h]
  \centering
  \caption{The four sections of the rubric and what each is designed to detect.}
  \label{tab:ef-mlts-sections}
  \begin{tabularx}{\linewidth}{@{}l X@{}}
    \toprule
    \textbf{Section (7 items each)} & \textbf{Failure mode it guards against} \\
    \midrule
    Data tests                  & Bad inputs reach training / serving (schema drift, missing features, costly features) \\
    Model development tests     & Models ship that are worse than a baseline, untuned, or untested on slices \\
    ML infrastructure tests     & Pipeline can't be reproduced, debugged, rolled back, or safely canary-deployed \\
    Monitoring tests             & Production failure goes undetected (silent drift, stale model, regressed accuracy) \\
    \bottomrule
  \end{tabularx}
\end{table}

One honesty rule is applied throughout: no point is awarded without written
evidence, and items that are genuinely not applicable to wholesale-price
forecasting (for example, demographic-fairness tests) are scored zero rather than
excused. That the rubric is self-administered is itself a threat to construct
validity; it is named as such in \S\ref{sec:ef-validity} and mitigated by the
recorded per-item evidence.

% -----------------------------------------------------------------------------
\section{Reproducibility tests}
\label{sec:ef-repro}

Reproducibility is measured in two complementary senses, following the
reproducibility-degrees framing of Gundersen and Kjensmo
\cite{gundersen2018reproducibility} and defined in Table~\ref{tab:ef-repro-tests}:
\emph{coverage}, a system-wide property asking whether every run carries enough
metadata to be reproduced, and \emph{demonstration}, a worked existence proof in
which one logged run is regenerated from scratch.

\begin{table}[h]
  \centering
  \caption{Two reproducibility tests.}
  \label{tab:ef-repro-tests}
  \begin{tabularx}{\linewidth}{@{}l X@{}}
    \toprule
    \textbf{Test} & \textbf{Definition} \\
    \midrule
    Provenance coverage   & \% of live runs in \texttt{algo.experiment\_runs} that carry the complete triple: \texttt{git\_sha}, (\texttt{data\_start\_utc}, \texttt{data\_end\_utc}), \texttt{model\_artifact\_id} \\
    Re-run demonstration  & Take one logged run; using only its committed config + target date, re-execute \texttt{LiveRunner}; report the max absolute difference vs the stored prediction \\
    \bottomrule
  \end{tabularx}
\end{table}

Both are reported because they answer different questions: coverage asks whether
the system \emph{could} reproduce any run, while the demonstration shows that it
\emph{does} for a concrete one. The deterministic LEAR configuration is used for
the demonstration because, given a fixed data window, it reproduces exactly;
the DNN and Transformer carry seed and ensemble stochasticity whose controlled
re-run is left to future work. Against the ACM Artifact Review and Badging
vocabulary~\cite{acm2020badging}, the artefact can be claimed Available,
Functional, and Reusable by self-review; the independent ``Reproduced'' badge,
which requires an external reviewer, is not claimed.

% -----------------------------------------------------------------------------
\section{Delivery throughput (DORA)}
\label{sec:ef-dora}

Delivery performance is reported using the throughput half of the DORA
metrics~\cite{forsgren2018accelerate}. One deployment is defined as one
successful run of \texttt{cd\_deploy.yml} on \texttt{main}, and lead time is the
wall-clock interval from the merge commit to the start of that run --- the
interval under the platform's automated control, rather than the time a feature
spends in development. Table~\ref{tab:ef-dora-scope} records which DORA metrics
are reported and which are deliberately withheld, with the reason in each case.

\begin{table}[h]
  \centering
  \caption{DORA: what is measured, what is excluded, and why.}
  \label{tab:ef-dora-scope}
  \begin{tabularx}{\linewidth}{@{}l c X@{}}
    \toprule
    \textbf{Metric} & \textbf{Reported?} & \textbf{Why / why not} \\
    \midrule
    Deployment frequency      & Yes & Mineable from git/CI history alone~\cite{forsgren2018accelerate} \\
    Lead time for changes     & Yes & Same; throughput metric \\
    Change failure rate       & \emph{No} & Requires an incident corpus the study window does not provide \\
    Mean time to restore      & \emph{No} & Same: needs labelled production incidents \\
    DORA tier placement       & \emph{No} & Tier thresholds are re-clustered yearly; unsafe to assign on a short window \\
    \bottomrule
  \end{tabularx}
\end{table}

% -----------------------------------------------------------------------------
\section{Drift monitoring counts}
\label{sec:ef-drift}

The drift test reports two counts over the study window, taken from
\texttt{algo.drift\_reports}: the number of drift reports written, and the number
flagged \texttt{drift\_detected = 1}. Deliberately, it does \emph{not} report a
drift-to-retrain latency. Under the reframing argued in
Chapter~\ref{ch:results-system} (\S\ref{sec:res-drift}), drift detection on this
platform serves monitoring and alerting rather than triggering, so a latency to
retrain is not a meaningful metric. That reframing is conditional: it holds while
retraining stays cheap, and would flip for expensive multi-year
retraining~\cite{mahadevan2024costaware} --- a boundary condition surfaced again
where the counts are reported.

% -----------------------------------------------------------------------------
\section{Maturity-model positioning}
\label{sec:ef-maturity}

The platform's per-practice maturity levels are reported in the qualitative
practice-by-practice table in the Results chapter
(Chapter~\ref{ch:results-system}, Table~\ref{tab:maturity}); this section
justifies the prior question of \emph{which} maturity model those levels are
assigned against. Zarour et al.~\cite{zarour2025} review five MLOps maturity
models (Table~\ref{tab:ef-maturity-candidates}), all descended from the
CMM/CMMI lineage. The platform is positioned against the Google three-level
model.

\begin{table}[h]
  \centering
  \caption{The five MLOps maturity models reviewed by Zarour et al.~\cite{zarour2025}, Table~7. All derive from CMM/CMMI.}
  \label{tab:ef-maturity-candidates}
  \begin{tabularx}{\linewidth}{@{}l c X@{}}
    \toprule
    \textbf{Model} & \textbf{Levels} & \textbf{Top level} \\
    \midrule
    Google                                  & 3 (0--2) & CI/CD pipeline automation \\
    Microsoft                               & 5 (0--4) & Full MLOps automated operations \\
    GigaOm                                  & 5 (0--4) & Comprehensive MLOps pipeline for frequent model updates \\
    John, Olsson \& Bosch (2021)            & 4 (0--3) & Fully-automated model monitoring \\
    Mateo-Casal\'i et al.\ (2022)            & 5 (0--4) & Optimization (pure CMMI mapping) \\
    \bottomrule
  \end{tabularx}
\end{table}

Google is chosen for three reasons. Its granularity matches the single-host MVP
scope: three levels avoid over-claiming distinctions --- such as those between
Microsoft's Levels 2--4 --- that a single deployment target cannot meaningfully
differentiate. Its top level, CI/CD-pipeline automation, is plausibly attainable
for this platform, whereas Microsoft's Level 4 (``full MLOps automated
operations'') would require multi-environment claims the platform does not make.
And it is an open, widely cited industry whitepaper~\cite{google_mlops_maturity},
the most-cited maturity reference in the Zarour review, giving the assignment an
external touchpoint a reviewer can check. The limitation is disclosed plainly:
all five models are descriptive ordinals that support \emph{positioning}, not
\emph{measurement}; the instruments of \S\ref{sec:ef-mlts}--\S\ref{sec:ef-drift}
carry the measurement load.

% -----------------------------------------------------------------------------
\section{Threats to validity (preview)}
\label{sec:ef-validity}

The detailed validity discussion accompanies the results in
\S\ref{sec:res-validity}; previewed here, with one line per threat type framed
per Runeson and H\"ost~\cite{runeson2009casestudy}, are the threats the
evaluation must defend against (Table~\ref{tab:ef-validity}).

\begin{table}[h]
  \centering
  \caption{Threats to validity at a glance; each is discussed under \S\ref{sec:res-validity}.}
  \label{tab:ef-validity}
  \begin{tabularx}{\linewidth}{@{}l X X@{}}
    \toprule
    \textbf{Threat} & \textbf{Source in this work} & \textbf{Mitigation} \\
    \midrule
    Construct       & ML Test Score is self-administered & Per-item evidence in \texttt{ml\_test\_score.md} \\
    Internal        & Single-system observation; drift conclusion conditional & State the boundary condition; objective measurements carry main load \\
    External        & Single-case design & Analytic generalisation only; no statistical claim about MLOps platforms \\
    Reliability     & Would another evaluator score the same? & Explicit scoring rules + recorded evidence per item \\
    \bottomrule
  \end{tabularx}
\end{table}

% -----------------------------------------------------------------------------
\section{What is not tested and why}
\label{sec:ef-non-claims}

Stating the evaluation's non-claims explicitly (Table~\ref{tab:ef-non-claims})
prevents the results from being over-read: each entry is something a reader might
reasonably expect but which the study window or design cannot support.

\begin{table}[h]
  \centering
  \caption{Explicit non-claims of this evaluation framework.}
  \label{tab:ef-non-claims}
  \begin{tabularx}{\linewidth}{@{}l X@{}}
    \toprule
    \textbf{Not claimed} & \textbf{Why} \\
    \midrule
    Long-run uptime / availability       & Prometheus tier is only days old; passive uptime over a short window is uninformative \\
    Change failure rate / MTTR           & No incident corpus in the window (see \S\ref{sec:ef-dora}) \\
    DORA tier placement                  & Tier thresholds re-clustered yearly; unsafe on a short window \\
    Comparison vs another MLOps platform & Single-case design; no controlled comparison \\
    Long-run forecasting accuracy        & Belongs to the modelling-side chapter (Ch.~\ref{ch:results-models}), not this framework \\
    Statistical generalisation           & Single case; only analytic generalisation per~\cite{runeson2009casestudy} \\
    \bottomrule
  \end{tabularx}
\end{table}
